% !TeX root = ../navier-stokes-manuscript.tex
% ============================================================
% 2.1 Vortex Stretching
% ============================================================

Axial stretching lengthens one narrow material vortex-tube segment; because its material volume stays fixed, its cross-section contracts, its rotating fluid moves inward, and the interior spins faster. We calculate this spin-up in the local axisymmetric case where the vorticity and stretching direction are aligned. The velocity field carrying the segment obeys
\[
  \underbrace{\partial_t u}_{\text{local acceleration}}
  +
  \underbrace{(u\cdot\nabla)u}_{\text{nonlinear transport}}
  =
  \underbrace{-\nabla p}_{\text{pressure}}
  +
  \underbrace{\nu\Delta u}_{\text{viscous diffusion}}.
\]
Taking the curl removes pressure and turns nonlinear transport into vorticity advection and stretching. The resulting rotational balance makes the segment's axial strain and changing cross-section the quantities that now need calculation.

\subsection{Vortex Stretching}\label{sub:ThePerfectlyStretchingVortex}

Let \(e\) be the aligned axis and \(L(t)\) the segment's length. Local axial strain is described by
\[
  S:=\frac12\bigl(\nabla u+(\nabla u)^{\mathsf T}\bigr),
  \qquad
  Se=\sigma(t)e,
  \qquad
  \sigma(t)>0,
  \qquad
  \frac{\dot L(t)}{L(t)}
  =
  e\cdot Se
  =
  \sigma(t).
\]
The positive strain lengthens the segment at logarithmic rate \(\sigma(t)\). Incompressibility keeps its material volume \(\mathcal V\) fixed, so its effective transverse area is \(A(t):=\mathcal V/L(t)\), and its effective radius is \(r_{\mathrm{eff}}(t):=\sqrt{A(t)/\pi}\). Thus
\[
  \nabla\cdot u=0,
  \qquad
  \frac{d}{dt}(A L)=0,
  \qquad
  \frac{\dot A}{A}=-\sigma(t),
  \qquad
  \frac{\dot r_{\mathrm{eff}}}{r_{\mathrm{eff}}}
  =
  -\frac{\sigma(t)}{2}.
\]
The logarithmic derivatives of area and radius are negative, while their positive contraction magnitudes are \(\sigma(t)\) and \(\sigma(t)/2\). As the same segment narrows, its aligned rotation is governed by \(\omega:=\nabla\times u\) and
\[
  \frac{D\omega}{Dt}
  =
  S\omega+\nu\Delta\omega,
  \qquad
  \frac{D}{Dt}
  :=
  \partial_t+u\cdot\nabla,
\]
and under the alignment \(\omega=|\omega|e\),
\[
  S\omega
  =
  \sigma(t)\omega,
  \qquad
  \frac12\frac{D|\omega|^2}{Dt}
  =
  \sigma(t)|\omega|^2
  +
  \nu\,\omega\cdot\Delta\omega.
\]
The segment's squared vorticity grows exactly where the complete displayed sum is positive. Aligned stretching contributes positively, while viscosity has the sign of \(\omega\cdot\Delta\omega\).

The energy identity proved in \Cref{prop:ClassicalKineticEnergyIdentity}, together with the antisymmetric part of \(\nabla u\), gives
\[
  \norm{u(t)}_2
  \leq
  \norm{u_0}_2
  <\infty,
  \qquad
  \left|\frac12\bigl(\nabla u-(\nabla u)^{\mathsf T}\bigr)\right|_{\mathrm F}
  =
  \frac{|\omega|}{\sqrt2}
  \leq
  |\nabla u|_{\mathrm F}.
\]
Finite kinetic energy does not prevent the rotational part of the velocity gradient from concentrating inside the narrowing segment. That failure of the global bound calls for a full-field comparison at a finer width.

\divider{\NavierStokes{} Scaling}

At a time \(t_0<T_*\), write \(\ell:=r_{\mathrm{eff}}(t_0)\). For \(\lambda>1\), define the exact \NavierStokes{} comparison at width \(\lambda^{-1}\ell\) by
\[
  u_\lambda(x,t)
  :=
  \lambda u(\lambda x,\lambda^2t),
  \qquad
  p_\lambda(x,t)
  :=
  \lambda^2p(\lambda x,\lambda^2t).
\]
At time \(\lambda^{-2}t_0\), the corresponding segment in \(u_\lambda\) has effective radius \(\lambda^{-1}\ell\). The rescaled field is another solution at another scale, not a later material state of the opening segment. Its momentum quantities transform as
\[
\begin{aligned}
  \partial_tu_\lambda(x,t)
  &=
  \lambda^3(\partial_tu)(\lambda x,\lambda^2t),
  \\
  (u_\lambda\cdot\nabla)u_\lambda(x,t)
  &=
  \lambda^3\bigl((u\cdot\nabla)u\bigr)(\lambda x,\lambda^2t),
  \\
  \nabla p_\lambda(x,t)
  &=
  \lambda^3(\nabla p)(\lambda x,\lambda^2t),
  \\
  \nu\Delta u_\lambda(x,t)
  &=
  \lambda^3\nu(\Delta u)(\lambda x,\lambda^2t).
\end{aligned}
\]
Every momentum term remains matched by the same cubic factor. Finer scale gives viscosity no relative advantage over nonlinear transport, even though Sobolev norms respond differently according to derivative order.

\divider{Critical Height}

Let \(\Lambda:=(-\Delta)^{1/2}\). The Fourier transform of the rescaled velocity satisfies
\[
  \widehat{u_\lambda}(\xi,t)
  =
  \lambda^{-2}\widehat u(\lambda^{-1}\xi,\lambda^2t),
\]
so a change of variables gives
\[
  \norm{u_\lambda(t)}_{\dot H^s}^2
  =
  \lambda^{2s-1}
  \norm{u(\lambda^2t)}_{\dot H^s}^2.
\]
The scaling power depends on derivative order: kinetic energy loses one power at \(s=0\), first derivatives gain one at \(s=1\), and the exponent vanishes at \(s=\tfrac12\). Name the half-order quantity by
\[
  H(t)
  :=
  \frac12\norm{u(t)}_{\dot H^{1/2}}^2
  =
  \frac12\norm{\Lambda^{1/2}u(t)}_2^2.
\]
Writing \(H_\lambda\) for the same quantity evaluated on \(u_\lambda\),
\[
  \norm{u_\lambda(t)}_2^2
  =
  \lambda^{-1}\norm{u(\lambda^2t)}_2^2,
  \qquad
  H_\lambda(t)=H(\lambda^2t),
  \qquad
  \norm{\nabla u_\lambda(t)}_2^2
  =
  \lambda\norm{\nabla u(\lambda^2t)}_2^2.
\]
The rescaled solution has less kinetic energy and a larger first-derivative norm, while its half-order height is scale neutral. That neutrality does not determine whether \(H\) rises or falls.

\divider{Nonlinear Production versus Viscous Dissipation}

At critical height, define the signed nonlinear production by
\[
  \mathcal P_{1/2}[u](t)
  :=
  -\operatorname{Re}
  \pair
    {\Lambda^{1/2}u(t)}
    {\Lambda^{1/2}\bigl((u(t)\cdot\nabla)u(t)\bigr)}_{L^2}.
\]
Pressure drops out of the pairing by incompressibility, while viscosity removes critical height at the nonnegative rate \(\nu\norm{\Lambda^{3/2}u}_2^2\). Consequently,
\[
  H'(t)
  +
  \nu\norm{\Lambda^{3/2}u(t)}_2^2
  =
  \mathcal P_{1/2}[u](t).
\]
The sign of \(H'(t)\) is set by nonlinear production relative to viscous dissipation. Under rescaling, both quantities change according to
\[
  \mathcal P_{1/2}[u_\lambda](t)
  =
  \lambda^2\mathcal P_{1/2}[u](\lambda^2t),
  \qquad
  \nu\norm{\Lambda^{3/2}u_\lambda(t)}_2^2
  =
  \lambda^2\nu\norm{\Lambda^{3/2}u(\lambda^2t)}_2^2.
\]
Neither gains ground through scale: both acquire the same quadratic factor. This scaling calculation has not determined the duration of a nonlinear transition between widths.

\divider{The Heat Clock}

The linear heat equation \(v_t=\nu\Delta v\) supplies a comparison through
\[
  \widehat v(\xi,t+\delta t)
  =
  e^{-\nu|\xi|^2\delta t}\widehat v(\xi,t).
\]
A scale-localized velocity structure of width \(r\) occupies frequencies \(|\xi|\simeq r^{-1}\). Its Fourier modes change by order one when
\[
  \nu|\xi|^2\delta t\simeq1,
\]
so the linear inverse-width comparison time is
\[
  \tau_\nu(r)
  :=
  \frac{r^2}{\nu}.
\]
Halving the width quarters this time, and a general reduction by \(\lambda\) gives
\[
  \tau_\nu(\lambda^{-1}r)
  =
  \lambda^{-2}\tau_\nu(r).
\]
This is a linear viscous comparison, not the realized duration of a nonlinear scale change.

\divider{The Zeno Cascade}

Choose dyadic widths first. For
\[
  r_n:=2^{-n}r_0,
  \qquad
  \tau_n:=\frac{r_n^2}{\nu},
\]
their heat comparison times shrink geometrically:
\[
  \tau_n
  =
  \frac{r_0^2}{\nu}4^{-n},
  \qquad
  \sum_{n=0}^{\infty}\tau_n
  =
  \frac{4r_0^2}{3\nu}
  <
  \infty.
\]
The convergence is finite arithmetic, not a realized schedule. One \NavierStokes{} velocity field would still have to exhibit localized widths
\[
  r_0>r_1>r_2>\cdots\longrightarrow0
\]
at sampled times
\[
  t_0<t_1<t_2<\cdots\uparrow T_*<\infty,
  \qquad
  t_{n+1}-t_n\asymp\tau_n,
\]
with comparison constants independent of \(n\). Its remaining time would then satisfy
\[
  T_*-t_n
  \asymp
  \sum_{k=n}^{\infty}\tau_k
  =
  \frac{4r_n^2}{3\nu}.
\]
At width \(r_n\), both the next sampled gap and the total time left are comparable to \(r_n^2/\nu\). Along this conditional history, the spatial cascade becomes a temporal demand on the same field because the scale-change intervals collapse as \(t\uparrow T_*\). Even this schedule does not control the field's first or successive Eulerian time derivatives at one fixed spatial point.
